\begin{abstractCN}
    服务器日志记录了系统运行过程中的事件状态、错误信息和组件行为，是开展故障诊断和异常检测的重要数据来源。针对原始日志格式不统一、动态参数较多、事件序列具有非规则时间间隔以及异常样本比例较低等问题，本文设计并实现了一种基于日志分析和 Liquid 架构的服务器异常检测系统。系统首先对原始日志进行字段抽取、时间排序、文本清洗和动态参数归一化处理；随后将非结构化日志正文解析为日志模板和事件 ID，并以时间窗口为单位构建事件频次、事件序列和相邻事件时间间隔等联合特征；最后使用 Liquid 模型对窗口级日志状态进行建模，并基于验证集选择异常判定阈值。实验部分以 BGL 数据集为主要测试对象，并引入 HDFS 数据集进行辅助对比。实验结果表明，本文系统能够完成从原始日志输入到异常结果输出的完整流程，在 BGL 测试集上取得 Precision 为 0.945、Recall 为 0.932、F1 值为 0.938 的检测结果。本文工作验证了将日志模板、统计频次、事件顺序和时间间隔信息结合用于服务器异常检测的可行性。

    \keywordsCN 日志分析；服务器异常检测；日志解析；特征工程；Liquid 架构；时序建模
\end{abstractCN}

\begin{abstractENG}
    Server logs record runtime events, error messages, and component behaviors, and therefore provide an important data source for fault diagnosis and anomaly detection. To address the problems of inconsistent log formats, numerous dynamic parameters, irregular time intervals between log events, and imbalanced anomaly samples, this thesis designs and implements a log-based server anomaly detection system using a Liquid architecture. The system first extracts key fields from raw logs, sorts log records by timestamp, cleans log messages, and normalizes dynamic parameters. It then parses unstructured log messages into log templates and event IDs, and constructs joint features based on event frequency, event sequence, and time interval within each time window. Finally, a Liquid model is used to model window-level log states, and the anomaly threshold is selected on the validation set. The experiments use the BGL dataset as the main benchmark and introduce the HDFS dataset for auxiliary comparison. The experimental results show that the proposed system can complete the end-to-end workflow from raw log input to anomaly output, achieving a Precision of 0.945, a Recall of 0.932, and an F1-score of 0.938 on the BGL test set. This work verifies the feasibility of combining log templates, statistical frequency, event order, and time-interval information for server anomaly detection.

    \keywordsENG Log Analysis; Server Anomaly Detection; Log Parsing; Feature Engineering; Liquid Architecture; Time Series Modeling
\end{abstractENG}
